\chapter{Stochastic Discount Factors and Equilibrium Pricing}

\begin{objectives}
Unify asset pricing through stochastic discount factors, derive risk premia and
no-arbitrage bounds, and connect equilibrium preferences to empirical moments.
\end{objectives}

\section{The fundamental pricing equation}

For payoff $X_{t+1}$ with price $P_t$, absence of arbitrage under suitable conditions
implies a positive stochastic discount factor $M_{t+1}$ such that
\[
P_t=\E_t[M_{t+1}X_{t+1}].
\]
For gross return $R_{i,t+1}=X_{i,t+1}/P_{i,t}$,
\[
1=\E_t[M_{t+1}R_{i,t+1}].
\]
This equation encompasses consumption-based equilibrium, state-price valuation, and
equivalent martingale measures.

The risk-free rate obeys $R_{f,t}=1/\E_t[M]$. Therefore
\[
\E_t[R_i]-R_f
=-\frac{\Cov_t(M,R_i)}{\E_t[M]}.
\]
Assets that pay when $M$ is high hedge bad states and require lower average returns.

\begin{theorem}[Hansen--Jagannathan bound]
For any excess return $R_i^e$ priced by $M$,
\[
\frac{\sigma_t(M)}{\E_t[M]}
\geq
\frac{\abs{\E_t[R_i^e]}}{\sigma_t(R_i^e)}.
\]
For a vector of returns, the strongest bound is the maximum conditional Sharpe ratio.
\end{theorem}
\begin{proof}
Pricing gives $\E[M R_i^e]=0$, so
$\E[R_i^e]=-\Cov(M,R_i^e)/\E[M]$. Apply Cauchy--Schwarz and divide by
$\sigma(R_i^e)$.
\end{proof}

The bound is diagnostic: a candidate consumption-based $M$ must be sufficiently
volatile to price observed Sharpe ratios. It does not identify the true SDF.

\section{Projection and the minimum-variance SDF}

Given gross returns $R$, any SDF can be decomposed into its projection on their span
plus an orthogonal component. The minimum-second-moment SDF in the payoff span is
\[
M^\star=a+b^\transpose R,
\]
with coefficients satisfying pricing equations. Unspanned components do not affect
prices of included assets but can matter for other claims.

In incomplete markets, many SDFs price the observed assets. Empirical asset pricing
therefore tests a joint hypothesis about the factor representation, conditioning
information, and set of test assets.

\section{Representative-agent equilibrium}

An agent maximizing $\E_t\sum_j\beta^ju(C_{t+j})$ has
\[
M_{t+1}=\beta\frac{u'(C_{t+1})}{u'(C_t)}.
\]
For lognormal consumption growth $g_{t+1}=\Delta\log C_{t+1}$ and CRRA utility,
\[
\log M_{t+1}=\log\beta-\gamma g_{t+1}.
\]
If $(g,r_i)$ are conditionally jointly normal, expected log excess return includes
consumption covariance and a Jensen term. Smooth aggregate consumption growth makes
the observed equity premium and low risk-free rate difficult to match with moderate
risk aversion, the equity-premium and risk-free-rate puzzles.

Extensions introduce habit, recursive utility, long-run risks, rare disasters,
heterogeneous agents, incomplete markets, or intermediary constraints. Each adds
state variables and identification demands; better fit alone does not establish the
mechanism.

\section{Conditional asset pricing}

If instruments $Z_t$ are known at $t$, conditional pricing implies unconditional
moments
\[
\E[(M_{t+1}R_{t+1}-\ones)\otimes Z_t]=0.
\]
Generalized method of moments estimates parameters by minimizing
$g_T(\theta)^\transpose W_Tg_T(\theta)$. Efficient weighting uses the inverse
long-run covariance of moments, but two-step estimation and weak identification
affect finite-sample inference.

The overidentifying-restrictions statistic tests whether all sample moments can be
close to zero. Rejection can reflect the SDF, return data, instruments, or maintained
regularity. Non-rejection does not prove economic truth.

\section{Entropy and good-deal restrictions}

No-arbitrage price intervals can be wide in incomplete markets. Good-deal bounds
restrict SDF volatility or admissible Sharpe ratios to exclude implausibly attractive
unspanned opportunities. Entropy methods choose a pricing measure closest to a
reference law subject to moment conditions:
\[
\min_{\Q} \E^\Prob\left[
\frac{\dd\Q}{\dd\Prob}\log\frac{\dd\Q}{\dd\Prob}
\right].
\]
This is an additional economic or statistical criterion, not a consequence of
no-arbitrage.

\section{Problems}

\begin{exercise}[Core: covariance pricing]
Derive expected excess return from the SDF equation and interpret the sign for an
insurance-like payoff.
\end{exercise}
\begin{exercise}[Core: HJ bound]
Derive the vector maximum-Sharpe form using the tangency portfolio.
\end{exercise}
\begin{exercise}[Advanced: GMM]
Write moments for a consumption-based model with CRRA utility, derive the Jacobian,
and describe HAC weighting with quarterly consumption and returns.
\end{exercise}


\chapter{CAPM, Multifactor Models, and Anomalies}

\begin{objectives}
Derive the CAPM, understand beta and factor-mimicking portfolios, evaluate
multifactor models cross-sectionally, and conduct anomaly research without common
biases.
\end{objectives}

\section{Mean--variance equilibrium}

Suppose investors have homogeneous beliefs over one-period means and covariance,
can borrow and lend at $R_f$, and choose portfolios by mean and variance. Every
investor holds a combination of the risk-free asset and the same tangency portfolio.
Market clearing makes that portfolio the market.

\begin{theorem}[CAPM]
Under the mean--variance equilibrium assumptions,
\[
\E[R_i]-R_f=\beta_i(\E[R_M]-R_f),\qquad
\beta_i=\frac{\Cov(R_i,R_M)}{\Var(R_M)}.
\]
\end{theorem}
\begin{proof}
The tangency portfolio first-order condition gives
$\mu-R_f\ones=\lambda\Sigma w_M$. Component $i$ is
$\mu_i-R_f=\lambda\Cov(R_i,R_M)$. Multiply by $w_M^\transpose$ and sum to obtain
$\mu_M-R_f=\lambda\Var(R_M)$; divide.
\end{proof}

The security market line is a statement about expected return and systematic
covariance, not total volatility. The zero-beta CAPM adapts the result when risk-free
borrowing is unavailable.

\section{Time-series beta estimation}

The market model
\[
R_{i,t}^e=\alpha_i+\beta_iR_{M,t}^e+\varepsilon_{i,t}
\]
estimates beta over a chosen window. Nonsynchronous trading, leverage change,
regime-dependent covariance, errors in the market proxy, and conditional betas
complicate interpretation. A statistically nonzero $\alpha$ is a joint rejection of
model, benchmark, implementation, and sampling assumptions.

\section{Arbitrage pricing theory}

If returns have approximate factor structure
$R_i=a_i+b_i^\transpose f+\varepsilon_i$ and idiosyncratic risk is diversifiable,
absence of arbitrage implies expected returns are approximately linear in factor
loadings:
\[
\E[R_i]=\lambda_0+b_i^\transpose\lambda.
\]
APT needs fewer preference restrictions than CAPM but does not name the factors.
Traded factors make $\lambda$ directly tied to factor expected returns; nontraded
macroeconomic factors need beta-pricing estimation.

\section{Empirical factor models}

Common factor families target size, value, profitability, investment, momentum,
quality, low risk, carry, and liquidity. A characteristic is an asset attribute used
to sort; a factor is a return. The distinction matters when a characteristic predicts
returns without representing covariance risk.

Portfolio construction choices include:

\begin{itemize}
\item universe and delisting returns;
\item breakpoint population and exchange;
\item weighting and rebalancing;
\item accounting-data publication lags;
\item microcap and liquidity filters;
\item neutralization and transaction costs.
\end{itemize}

Published factor data are excellent benchmarks but do not replace reconstruction
when the research question concerns implementability.

\section{Cross-sectional tests}

The Fama--MacBeth procedure estimates cross-sectional regressions each period:
\[
R_{i,t}=a_t+\beta_i^\transpose\lambda_t+u_{i,t},
\]
then averages $\widehat\lambda_t$. Time-series standard errors of the premia must
handle autocorrelation. Estimated betas introduce errors-in-variables; Shanken-type
corrections address one component under strong assumptions.

Testing factors on portfolios sorted by the same characteristic can make power and
interpretation circular. Include diverse test assets, report pricing errors in
economic units, and compare against nested and nonnested benchmarks.

\section{Anomalies and the research process}

An anomaly may reflect compensation for risk, behavioral mispricing, institutional
friction, data mining, or implementation error. A credible study:

\begin{enumerate}
\item states the economic mechanism before examining the test sample;
\item freezes feature definitions and availability dates;
\item accounts for delistings, survivorship, and restatements;
\item uses an out-of-sample or independent-market test;
\item reports turnover, spread, impact, borrow, and capacity;
\item compares to simple factor exposures;
\item publishes failures and sensitivity.
\end{enumerate}

The probability of an extreme backtest rises with the number and dependence of tried
strategies. Deflated Sharpe ratios and reality-check procedures attempt selection
adjustments but require an honest count or model of the search process.

\section{Problems}

\begin{exercise}[Core: CAPM]
Derive the zero-beta version when there is no risk-free asset but all investors share
the same efficient frontier.
\end{exercise}
\begin{exercise}[Core: factor neutrality]
Given loading matrix $B$, formulate the closest portfolio to target $w_0$ subject to
$B^\transpose w=0$ and full investment.
\end{exercise}
\begin{exercise}[Advanced: two-pass bias]
Explain how noisy first-pass beta estimates bias the second-pass risk-premium
regression. Propose portfolios or instruments that improve precision.
\end{exercise}


\chapter{Return Predictability and Present-Value Relations}

\begin{objectives}
Derive present-value decompositions, connect valuation ratios to cash-flow and
discount-rate news, and evaluate predictive regressions with appropriate skepticism.
\end{objectives}

\section{Log-linear present value}

Let $p_t$ and $d_t$ be log price and dividend. Log return
$r_{t+1}=\log(P_{t+1}+D_{t+1})-\log P_t$ can be approximated around the average
dividend-price ratio:
\[
r_{t+1}\approx k+\rho p_{t+1}+(1-\rho)d_{t+1}-p_t.
\]
Rearranging and iterating under a terminal condition yields
\[
d_t-p_t
\approx -\frac{k}{1-\rho}
+\sum_{j=0}^{\infty}\rho^j
\left(\E_t r_{t+1+j}-\E_t\Delta d_{t+1+j}\right).
\]
A high dividend-price ratio must forecast high future returns, low future dividend
growth, or both within the approximation.

\section{News decomposition}

Unexpected return can be decomposed into cash-flow news and discount-rate news.
Prices can fall on good cash-flow news if discount rates rise more. Interpreting a
price response without separating the two is incomplete.

VAR-based decompositions estimate revisions to long-horizon expectations, so results
depend on included state variables, persistence, sample, and terminal assumptions.

\section{Predictive regressions}

A standard regression is
\[
r_{t+1}=\alpha+\beta x_t+\varepsilon_{t+1}.
\]
When $x_t$ is persistent and its innovation correlates with $\varepsilon_{t+1}$,
finite-sample bias and nonstandard behavior arise. Overlapping long-horizon returns
create moving-average residuals and mechanically smooth noise.

Out-of-sample $R^2$ relative to historical mean is
\[
R^2_{\mathrm{OS}}=
1-\frac{\sum_t(r_t-\widehat r_t)^2}
{\sum_t(r_t-\bar r_{t-1})^2}.
\]
Small positive values can be economically meaningful for volatile returns, but
statistical comparison must reflect real-time estimation and model selection.

\section{Economic value}

A mean--variance investor might allocate
\[
w_t=\frac{\widehat\mu_t}{\gamma\widehat\sigma_t^2},
\]
subject to leverage and turnover constraints. Certainty-equivalent improvement
evaluates forecasts through a decision. This exercise is sensitive to scaling,
variance forecast, parameter caps, and costs.

Forecast combinations reduce model-selection risk. Simple equal weighting is a
strong benchmark; optimal combination weights are themselves estimated and can be
unstable.

\section{Behavioral and institutional mechanisms}

Underreaction can arise from slow information diffusion or anchoring; overreaction
from extrapolation or constrained arbitrage. Limits to arbitrage include short-sale
risk, funding, horizon mismatch, and noise-trader risk. A behavioral narrative must
make distinct predictions about horizons, investor types, events, or limits; otherwise
it is an after-the-fact label.

\section{Problems}

\begin{exercise}[Core: present value]
Complete the iteration of the log-linear return identity and state the no-bubble
terminal condition.
\end{exercise}
\begin{exercise}[Core: overlapping returns]
Show that a $k$-period overlapping return induces an MA$(k-1)$ structure even when
one-period returns are iid.
\end{exercise}
\begin{exercise}[Research]
Compare valuation-ratio forecasts against the historical mean using recursive
estimation. Report statistical loss, certainty equivalent, turnover, and performance
across predeclared subperiods.
\end{exercise}
